
This section formalizes the protocol's state, parameters, and mathematical invariants. All subsequent mechanism definitions rely on the notation and constants introduced here. We adopt a fixed-point arithmetic model (WAD scaling, $10^{18}$) to ensure precision across all computations.

\section{State Variables}

The protocol state is partitioned into market state, time state, and index state. Each variable has a well-defined domain that is enforced by construction.

\subsection{Market State}

\begin{center}
\begin{tabular}{llll}
\toprule
\textbf{Symbol} & \textbf{Type} & \textbf{Domain} & \textbf{Description} \\
\midrule
$S_{SY}$ & u256 & $[0, 2^{256})$ & SY reserve in market \\
$S_{PT}$ & u256 & $[0, 2^{256})$ & PT reserve in market \\
$L$ & u256 & $[0, 2^{256})$ & Total LP token supply \\
$r$ & u256 & $[0, 4.6 \cdot 10^{18}]$ & ln(implied rate) in WAD \\
$r_{root}$ & u256 & $[0, 10^{20}]$ & Scalar root parameter \\
$r_{anchor}$ & u256 & $[0, 10^{19}]$ & Rate anchor parameter \\
$\phi$ & u256 & $[0, 0.05 \cdot 10^{18}]$ & Fee rate in WAD \\
\bottomrule
\end{tabular}
\end{center}

\textbf{Assumption}: $S_{PT} + S_{SY} > 0$ at all times due to enforced minimum liquidity lock.

\subsection{Time State}

\begin{center}
\begin{tabular}{llll}
\toprule
\textbf{Symbol} & \textbf{Type} & \textbf{Domain} & \textbf{Description} \\
\midrule
$T$ & u64 & Unix timestamp & Expiry timestamp \\
$t$ & u64 & Unix timestamp & Current block timestamp \\
$\tau$ & u64 & $[0, T]$ & Time to expiry (seconds) \\
$\tau_y$ & u256 & $[0, \infty)$ & Time to expiry (years, WAD-scaled) \\
\bottomrule
\end{tabular}
\end{center}

\subsection{Index State}

\begin{center}
\begin{tabular}{llll}
\toprule
\textbf{Symbol} & \textbf{Type} & \textbf{Domain} & \textbf{Description} \\
\midrule
$I_{global}$ & u256 & $[10^{18}, \infty)$ & Global PY index (watermark) \\
$I_{user}$ & u256 & $[10^{18}, \infty)$ & User's last recorded index \\
\bottomrule
\end{tabular}
\end{center}

\section{Constants}

The protocol relies on the following global constants, which are fixed for the lifetime of a given market instance.

\textbf{Scaling Constants:}
\begin{itemize}
\item \textbf{WAD} ($W = 10^{18}$): Fixed-point scaling factor for all amounts and rates
\item \textbf{Seconds per year} ($Y = 31{,}536{,}000$): Time normalization constant for annualized rates
\end{itemize}

\textbf{Safety Bounds:}
\begin{itemize}
\item \textbf{Minimum liquidity} ($L_{min} = 1000$): Locked LP to prevent division by zero
\item \textbf{Max ln implied rate} ($r_{max} = 4.6 \cdot 10^{18}$): Caps implied rate at approximately 10,000\% APY
\item \textbf{Max proportion} ($p_{max} = 0.999 \cdot 10^{18}$): Prevents extreme pool imbalance
\item \textbf{Min proportion} ($p_{min} = 0.001 \cdot 10^{18}$): Prevents extreme pool imbalance
\end{itemize}

\section{Fixed-Point Arithmetic}

The protocol uses a dual-layer precision model to balance gas efficiency with computational accuracy.

\subsection{WAD Arithmetic (Interface Layer)}

All token amounts and external-facing values use WAD scaling:

\[
\operatorname{wad\_mul}(a, b) = \left\lfloor \frac{a \times b + W/2}{W} \right\rfloor
\]

\[
\operatorname{wad\_div}(a, b) = \left\lfloor \frac{a \times W + b/2}{b} \right\rfloor
\]

where $W = 10^{18}$ and rounding is to nearest integer.

All intermediate multiplications are performed using widened precision to prevent overflow. Rounding is symmetric and biased toward the nearest integer to minimize cumulative drift.

\subsection{cubit 64.64 (Internal Layer)}

Transcendental functions use 64.64 binary fixed-point representation with 64 integer bits and 64 fractional bits, providing approximately 19 decimal digits of precision.

The 64.64 format represents a number $x$ as:
\[
x_{\text{fixed}} = x \times 2^{64}
\]

\paragraph{Conversion Functions.}

\[
\operatorname{wad\_to\_fp}(x) = \left\lfloor \frac{x \times 2^{64}}{10^{18}} \right\rfloor
\]

\[
\operatorname{fp\_to\_wad}(x) = \left\lfloor \frac{x \times 10^{18}}{2^{64}} \right\rfloor
\]

\paragraph{High-Precision Functions.}

The library provides the following transcendental functions:

\begin{itemize}
\item $\operatorname{exp\_wad}(x)$: Exponential function, takes WAD input and returns WAD output
\item $\operatorname{exp\_neg\_wad}(x)$: Negative exponential for decay functions
\item $\operatorname{ln\_wad}(x)$: Natural logarithm, returns (value, is\_negative) tuple
\item $\operatorname{pow\_wad}(x, y)$: Power function computing $x^y$
\item $\operatorname{sqrt\_wad}(x)$: Square root function, WAD-normalized output
\end{itemize}

All transcendental functions clamp outputs on domain violation and return bounded values to preserve protocol safety.

\section{Time Model}

Time is measured in Unix timestamps (seconds since epoch):

\[
\tau = T - t
\]

where $T$ is expiry timestamp and $t$ is current block timestamp.

Time expressed in years (WAD-scaled):

\[
\tau_y = \frac{\tau \times W}{Y}
\]

For $t \geq T$, the protocol treats $\tau = 0$ and transitions to post-expiry redemption semantics.

\section{Core Equations}

This section defines the mathematical relationships governing AMM pricing, yield computation, and fee adjustment.

\subsection{Pool Proportion and Logit Transformation}

The proportion of PT in the pool determines market pricing:

\[
p = \frac{S_{PT}}{S_{PT} + S_{SY}}
\]

\textbf{Domain constraint}: $p_{min} \leq p \leq p_{max}$. Trades that would violate this constraint are rejected.

To map bounded pool proportions to an unbounded rate domain, the protocol uses a logit transformation:

\[
\operatorname{logit}(p) = \ln\left(\frac{p}{1-p}\right)
\]

Properties:
\begin{itemize}
    \item $\operatorname{logit}(0.5) = 0$
    \item $\operatorname{logit}(p) \to -\infty$ as $p \to 0$
    \item $\operatorname{logit}(p) \to +\infty$ as $p \to 1$
\end{itemize}

\subsection{Pricing Model}

Time-to-expiry is incorporated through a scalar that dampens price sensitivity as maturity approaches.

\paragraph{Rate Scalar.}

\[
r_{scalar} = \frac{r_{root} \times Y}{\tau}
\]

As $\tau \to 0$: $r_{scalar} \to \infty$, making prices insensitive to pool balance changes.

\paragraph{Exchange Rate.}

The PT/SY exchange rate derived from market state:

\[
\operatorname{exchangeRate} = \frac{\operatorname{logit}(p)}{r_{scalar}} + r_{anchor}
\]

\paragraph{Implied Rate.}

The annualized implied yield from the exchange rate:

\[
r = \ln(\operatorname{exchangeRate}) \times \frac{Y}{\tau}
\]

\textbf{Constraint}: $0 \leq r \leq r_{max}$. Values outside this range are clamped.

\paragraph{PT Price.}

PT price in terms of SY, derived from implied rate:

\[
P_{PT} = e^{-r \times \tau_y / W}
\]

where $\tau_y$ is time to expiry in WAD-scaled years.

Properties:
\begin{itemize}
    \item $P_{PT} < 1$ before expiry (trades at discount)
    \item $P_{PT} \to 1$ as $\tau \to 0$ (converges at maturity)
\end{itemize}

\paragraph{YT Price.}

YT price is derived directly from the conservation invariant:

\[
P_{YT} = 1 - P_{PT}
\]

Properties:
\begin{itemize}
    \item $P_{YT} > 0$ before expiry
    \item $P_{YT} \to 0$ as $\tau \to 0$ (expires worthless)
\end{itemize}

\subsection{Rate Anchor Dynamics}

The rate anchor maintains implied rate continuity across trades, preventing discontinuous price jumps.

\paragraph{Anchor Calculation.}

Given the last stored implied rate $r_{last}$ and current market state:

\[
\operatorname{exchangeRate}_{target} = e^{r_{last} \times \tau / Y}
\]

\[
r_{anchor} = \operatorname{exchangeRate}_{target} - \frac{\operatorname{logit}(p)}{r_{scalar}}
\]

\paragraph{Anchor Update.}

After each trade, the anchor is recalculated from the new implied rate to ensure smooth price transitions.

\subsection{Fee Model}

Swap fees decay exponentially toward expiry, reducing friction as price convergence becomes more certain.

\[
\phi_{adjusted} = e^{\ln(\phi_{root}) \times \tau / Y}
\]

where $\phi_{root}$ is the base fee rate (e.g., 0.001 WAD = 0.1\%).

\textbf{Requirement}: $0 < \phi_{root} < 1$ to ensure monotonic decay of swap fees.

\begin{center}
\begin{tabular}{ll}
\toprule
\textbf{Time to Expiry} & \textbf{Adjusted Fee (if base = 0.1\%)} \\
\midrule
1 year & 0.100\% \\
6 months & $\approx 0.070\%$ \\
1 month & $\approx 0.022\%$ \\
1 week & $\approx 0.006\%$ \\
At expiry & $\approx 0\%$ \\
\bottomrule
\end{tabular}
\end{center}

\subsection{Yield Index Model}

The PY index accounts for cumulative yield attributable to YT holders using a monotonic watermark mechanism.

\paragraph{Index Update.}

\[
I_{global} = \max(I_{global}, I_{current})
\]

where $I_{current}$ is the current exchange rate from the oracle.

The watermark ensures $I_{global}$ is monotonically non-decreasing. This prevents yield accounting errors during temporary rate decreases.

\paragraph{Interest Calculation.}

For a user with YT balance $B_{YT}$ and last recorded index $I_{user}$:

\[
\text{Interest} = B_{YT} \times \frac{I_{global} - I_{user}}{I_{user}}
\]

After claiming, the user's index is updated: $I_{user} \leftarrow I_{global}$

\section{Invariants}

The following invariants are safety-critical. Any violation constitutes a protocol failure and must be prevented by construction or rejected at runtime.

\begin{center}
\begin{tabularx}{\textwidth}{l >{\raggedright\arraybackslash}X l}
\toprule
\textbf{Invariant} & \textbf{Expression} & \textbf{Enforcement} \\
\midrule
Supply equality & $\text{PT.supply} = \text{YT.supply}$ & Contract logic \\
Conservation & PT + YT = SY (value) & Arbitrage \\
Index watermark & $I_{global}^{t+1} \geq I_{global}^{t}$ & Contract logic \\
Proportion bounds & $p_{min} \leq p \leq p_{max}$ & Trade rejection \\
Rate bounds & $0 \leq r \leq r_{max}$ & Clamping \\
Reserve positivity & $S_{SY} > 0 \land S_{PT} > 0$ & LP minimum lock \\
\bottomrule
\end{tabularx}
\end{center}
